% Constructors: building blocks that take the data of a mathematical object
% and draw it with the dzg TikZ library vocabulary. The objects themselves are
% figures, one file each, under figures/objects/<family>/, assembled from these
% constructors; a figure places an object with \pic (<name>) {object=...} and
% decorates it.
% \input by tikzlibrarydzg.code.tex, so "@" is a letter here.

% \pic (Q) at (x,y) [<keys>] {object=<family>/<name>}: the object in
% figures/objects/<family>/<name>.tikz. Its points, named -<point> in the
% file, are Q-<point>; keys such as `ade labels=none` or `ell={...}` reach it.
% The picture's scale spaces the object's coordinates (transform shape) and
% leaves its node sizes absolute (transform shape=false inside).
\tikzset{pics/object/.style={/tikz/transform shape, code={%
  \tikzset{transform shape=false}\input{objects/#1.tikz}}}}

% =============================================================================
% Lattice polygons: constructors
% =============================================================================
% Building blocks for lattice polygons, Symington polygons, integral-affine
% singularities and toric fans. A constructor takes the data of one polygon
% and draws it; the polygons themselves are objects, one file each, under
% figures/objects/ (ade/, ias/, toric/), assembled from these constructors.
% \input by dzg-constructors.tex, so "@" is a letter here.
%
% Commands:
%   \adepolygon{<px>/<py>}{<x/y, ...>}
%       the toric ADE pair (Y, C) of the lattice polygon Q with distinguished
%       point p* and vertices v_1, ..., v_k after p* counterclockwise
%       (Alexeev-Thompson, arXiv:1712.07932, Notation 3.7); see below.
%   \adetitle{<math>}           the decorated Dynkin symbol under Q
%   \symingtonpolygon[<keys>]{<sides>}{<surgeries>}
%       a polygon given by edge vectors l_i d_i, with Symington surgeries;
%       see below.
%   \setell{<first index>}{<l_first, l_first+1, ...>}
%       the values l_i the polygon is drawn at, unless the figure passes
%       `ell={...}`.
%   \iassingularity{<point>}{<charge>}   an integral-affine singular point
%   \fanrays{<i/a/b/anchor, ...>}       rays rho_i = (a,b) of a complete fan
%   \fancones{<i/coordinate, ...>}      the maximal cones sigma_i
%   \dzgmonomial{<i>}{<j>}              x^i y^j
%
% Keys:
%   ade labels=monomials|none   monomials, side lengths, p* (default monomials)
%   ade title=true|false        the decorated symbol (default true)
%   ade affine                  p* is interior to the side [v_k, v_1]
%   ias labels=canonical|multiplicities|none
%                               canonical: point names, edge vectors, charges;
%                               multiplicities: charges alone
%   ias involutions             draw the involutions an object records
%   ell={l_first, ...}          the values l_i of a Symington polygon
%   fan labels=canonical|none, fan length=<factor>
%
% Conditionals for objects: \ifiaslabels, \ifiasmultiplicities,
% \ifiasinvolutions, \ifadelabels, \ifadetitle, \ifadeaffine, \iffanlabels.

\newif\ifiaslabels \iaslabelstrue
\newif\ifiasmultiplicities \iasmultiplicitiestrue
\newif\ifiasinvolutions
\newif\ifadelabels \adelabelstrue
\newif\ifadetitle \adetitletrue
\newif\ifadeaffine
\newif\iffanlabels \fanlabelstrue
\def\fanlength{2}
\let\dzg@ell@given\relax
\tikzset{
  ias labels/.is choice,
  ias labels/canonical/.code={\iaslabelstrue\iasmultiplicitiestrue},
  ias labels/none/.code={\iaslabelsfalse\iasmultiplicitiesfalse},
  ias labels/multiplicities/.code={\iaslabelsfalse\iasmultiplicitiestrue},
  ias involutions/.is if=iasinvolutions,
  % text of an edge vector or point name of an IAS
  ias label/.style={font=\footnotesize, inner sep=2pt},
  ade labels/.is choice,
  ade labels/monomials/.code={\adelabelstrue},
  ade labels/none/.code={\adelabelsfalse},
  ade title/.is if=adetitle,
  ade affine/.is if=adeaffine,
  ell/.code={\def\dzg@ell@given{#1}},
  fan labels/.is choice,
  fan labels/canonical/.code={\fanlabelstrue},
  fan labels/none/.code={\fanlabelsfalse},
  fan length/.store in=\fanlength,
}

% x^i y^j, written 1, x, x^2, y, xy, ...
\newcommand{\dzgmonomial}[2]{%
  \ifnum#1=0 \ifnum#2=0 1\fi\else x\ifnum#1>1 ^{#1}\fi\fi
  \ifnum#2>0 y\ifnum#2>1 ^{#2}\fi\fi}

% \iassingularity{<point>}{<charge>}: the charge is printed when it is not 1
% and multiplicities are shown.
\newcommand\iassingularity[2]{%
  \ifiasmultiplicities\ifnum#2>1
    \node[ias singularity=#2] at (#1) {};%
  \else\node[ias singularity] at (#1) {};\fi
  \else\node[ias singularity] at (#1) {};\fi}

% =============================================================================
% Symington polygons
% =============================================================================


% Source of the construction: the (18,2,0)_1 and (18,0,0)_1 templates of the
% dissertation, research/writing/dissertation/sections/2-part-moduli/
% 5-chapter-ias-dlt/500-ias.md (B(lambda) = P(lambda) u P(lambda)^op), after
% Symington 2003 and Alexeev-Engel-Garza-Schaffler 2023 §4.2. The barycentric
% coordinates l_i = lambda . alpha_i of the Vinberg roots alpha_i fix the
% sides: side i is
%   v_i = l_i d_i            (kind w),
%   v_i = (l_i / 2) d_i      (kind h),
%   v_i = (l_i + l_j / 2) d_i  (kind s: a Symington surgery with root j),
% with d_i the primitive direction of the side, counterclockwise. The surgery
% cuts the segment of lattice length l_j / 2 starting at distance l_i / 2 along
% the side and folds it to the singular point
%   apex = (cut start or cut end) + (l_j / 2) u_j,
% so the cut triangle is unimodular up to scale.
%
% Data:
%   sides      <i>/<dx>/<dy>/<kind>/<j>   (j = 0 when the side has no surgery)
%   surgeries  <i>/<j>/<ux>/<uy>/<anchor: a = cut start, b = cut end>/<apex>
% Points: -p<i> the vertices, -p<apex> the singular points, -c<j>a and -c<j>b
% the ends of the cut of surgery j.
% Keys of \symingtonpolygon: first=<index of the first vertex>, last=<index of
% the last vertex>, closed (the last side returns to p_first), charges={i/c,
% ...}. Every vertex is an integral-affine singular point of charge 1 unless
% `charges` says otherwise; vertices joined by sides of length 0 add up.
% \dzg@ias@ell{<i>}: the value l_i the polygon is drawn at.
\def\dzg@ias@ell#1{\csname dzg@ias@l@#1\endcsname}
% \dzg@ias@setell{<first index>}{<l list>}
\def\dzg@ias@setell#1#2{%
  \foreach \dzg@v [count=\dzg@c from #1] in {#2}
    {\expandafter\xdef\csname dzg@ias@l@\dzg@c\endcsname{\dzg@v}}}

% \dzg@ias@sidelen{<i>}{<kind>}{<j>}: sets \dzg@len, the lattice length.
\def\dzg@ias@sidelen#1#2#3{%
  \if#2w\pgfmathsetmacro\dzg@len{\dzg@ias@ell{#1}}\fi
  \if#2h\pgfmathsetmacro\dzg@len{\dzg@ias@ell{#1}/2}\fi
  \if#2s\pgfmathsetmacro\dzg@len{\dzg@ias@ell{#1}+\dzg@ias@ell{#3}/2}\fi}

% \dzg@ias@vector{<i>}{<kind>}{<j>}{<dx>}{<dy>}: the edge vector v_i.
\def\dzg@ias@vector#1#2#3#4#5{%
  \if#2w$v_{#1}=\ell_{#1}\cdot(#4,#5)$\fi
  \if#2h$v_{#1}=\tfrac12\ell_{#1}\cdot(#4,#5)$\fi
  \if#2s$v_{#1}=(\ell_{#1}+\tfrac12\ell_{#3})\cdot(#4,#5)$\fi}

% \dzg@ias@symington{<first>}{<last vertex>}{<closed: 1|0>}{<sides>}{<surgeries>}
%   {<charges i/c: the charge of vertex p_i where it is not 1>}
\def\dzg@ias@symington#1#2#3#4#5#6{%
  % Vertices, as numbers \dzg@ias@x@<i>, \dzg@ias@y@<i> and coordinates -p<i>.
  \expandafter\xdef\csname dzg@ias@x@#1\endcsname{0}%
  \expandafter\xdef\csname dzg@ias@y@#1\endcsname{0}%
  \foreach \i/\dx/\dy/\kind/\j in #4 {%
    \dzg@ias@sidelen{\i}{\kind}{\j}%
    \expandafter\xdef\csname dzg@ias@len@\i\endcsname{\dzg@len}%
    \pgfmathtruncatemacro\dzg@next{\i+1}%
    \pgfmathsetmacro\dzg@x{\csname dzg@ias@x@\i\endcsname+\dzg@len*\dx}%
    \pgfmathsetmacro\dzg@y{\csname dzg@ias@y@\i\endcsname+\dzg@len*\dy}%
    \expandafter\xdef\csname dzg@ias@x@\dzg@next\endcsname{\dzg@x}%
    \expandafter\xdef\csname dzg@ias@y@\dzg@next\endcsname{\dzg@y}}%
  \pgfmathtruncatemacro\dzg@ias@n{#2-#1+1}%
  \xdef\dzg@ias@cx{0}\xdef\dzg@ias@cy{0}%
  \foreach \i in {#1,...,#2} {%
    \coordinate (-p\i) at (\csname dzg@ias@x@\i\endcsname,\csname dzg@ias@y@\i\endcsname);
    \pgfmathsetmacro\dzg@x{\dzg@ias@cx+\csname dzg@ias@x@\i\endcsname/\dzg@ias@n}%
    \pgfmathsetmacro\dzg@y{\dzg@ias@cy+\csname dzg@ias@y@\i\endcsname/\dzg@ias@n}%
    \xdef\dzg@ias@cx{\dzg@x}\xdef\dzg@ias@cy{\dzg@y}}%
  % Surgeries: cut ends -c<j>a, -c<j>b and the singular point -p<apex>.
  \foreach \i/\j/\ux/\uy/\anchor/\apex in #5 {%
    \pgfmathtruncatemacro\dzg@next{\i+1}%
    \ifnum\dzg@next>#2 \def\dzg@next{#1}\fi
    \edef\dzg@px{\csname dzg@ias@x@\i\endcsname}\edef\dzg@py{\csname dzg@ias@y@\i\endcsname}%
    \edef\dzg@qx{\csname dzg@ias@x@\dzg@next\endcsname}%
    \edef\dzg@qy{\csname dzg@ias@y@\dzg@next\endcsname}%
    \pgfmathsetmacro\dzg@s{\dzg@ias@ell{\i}/2/max(\csname dzg@ias@len@\i\endcsname,0.001)}%
    \pgfmathsetmacro\dzg@t{(\dzg@ias@ell{\i}+\dzg@ias@ell{\j})/2/max(\csname dzg@ias@len@\i\endcsname,0.001)}%
    % Direction of the side, the cut ends and the apex, as numbers.
    \pgfmathsetmacro\dzg@dx{(\dzg@qx-\dzg@px)}\pgfmathsetmacro\dzg@dy{(\dzg@qy-\dzg@py)}%
    \pgfmathsetmacro\dzg@ax{\dzg@px+\dzg@s*\dzg@dx}\pgfmathsetmacro\dzg@ay{\dzg@py+\dzg@s*\dzg@dy}%
    \pgfmathsetmacro\dzg@bx{\dzg@px+\dzg@t*\dzg@dx}\pgfmathsetmacro\dzg@by{\dzg@py+\dzg@t*\dzg@dy}%
    \if\anchor b\let\dzg@ox\dzg@bx\let\dzg@oy\dzg@by
    \else\let\dzg@ox\dzg@ax\let\dzg@oy\dzg@ay\fi
    \pgfmathsetmacro\dzg@x{\dzg@ox+\dzg@ias@ell{\j}/2*\ux}%
    \pgfmathsetmacro\dzg@y{\dzg@oy+\dzg@ias@ell{\j}/2*\uy}%
    \expandafter\xdef\csname dzg@ias@ax@\apex\endcsname{\dzg@x}%
    \expandafter\xdef\csname dzg@ias@ay@\apex\endcsname{\dzg@y}%
    \coordinate (-c\j a) at (\dzg@ax,\dzg@ay);
    \coordinate (-c\j b) at (\dzg@bx,\dzg@by);
    \coordinate (-p\apex) at (\dzg@x,\dzg@y);}%
  \begin{scope}[on background layer]
    \fill[ias region] (-p#1) \foreach \i in {#1,...,#2} { -- (-p\i) } -- cycle;
    \foreach \i/\j/\ux/\uy/\anchor/\apex in #5 {%
      \ifdim\dzg@ias@ell{\j}pt>0pt
        \path[ias surgery] (-c\j a) -- (-p\apex) -- (-c\j b) -- cycle;\fi}
  \end{scope}
  % Sides, oriented counterclockwise; a surgery side is cut in three.
  \foreach \i/\dx/\dy/\kind/\j in #4 {%
    \pgfmathtruncatemacro\dzg@next{\i+1}%
    \ifnum\dzg@next>#2 \def\dzg@next{#1}\fi
    \ifdim\csname dzg@ias@len@\i\endcsname pt>0pt
      \ifnum\j=0
        \draw[fan ray] (-p\i) -- (-p\dzg@next);
      \else
        \ifdim\dzg@ias@ell{\i}pt>0pt \draw[fan ray, -] (-p\i) -- (-c\j a);\fi
        \ifdim\dzg@ias@ell{\j}pt>0pt
          \draw[fan ray, -, dashed, draw opacity=0.4] (-c\j a) -- (-c\j b);\fi
        \ifdim\dzg@ias@ell{\i}pt>0pt \draw[fan ray] (-c\j b) -- (-p\dzg@next);\fi
      \fi
      \ifiaslabels
        \pgfmathsetmacro\dzg@a{atan2(-\dx,\dy)+180}%
        \path ($(-p\i)!0.5!(-p\dzg@next)$)
          node[ias label, anchor=\dzg@a] {\dzg@ias@vector{\i}{\kind}{\j}{\dx}{\dy}};
      \fi
    \fi}%
  % Surgery vectors v_j and their singular points.
  \foreach \i/\j/\ux/\uy/\anchor/\apex in #5 {%
    \ifdim\dzg@ias@ell{\j}pt>0pt
      \draw[fan ray, draw=dzg singular] (-c\j\anchor) -- (-p\apex);
      \ifiaslabels
        % Past the singular point, in the direction of the vector.
        \pgfmathsetmacro\dzg@a{atan2(-\uy,-\ux)}%
        \path (-p\apex) node[ias label, text=dzg singular, anchor=\dzg@a,
          inner sep=6pt] {$v_{\j}=\tfrac12\ell_{\j}\cdot(\ux,\uy)$};
      \fi
    \fi}%
  % Singular points: each apex with the number of apexes at the same point,
  % each vertex with the number of vertices at the same point (sides of length
  % 0 collapse vertices) plus its extra multiplicity.
  \foreach \i/\j/\ux/\uy/\anchor/\apex in #5 {%
    \ifdim\dzg@ias@ell{\j}pt>0pt
      \gdef\dzg@mult{0}\gdef\dzg@first{1}%
      \foreach \k/\l/\vx/\vy/\b/\bapex in #5 {%
        \ifdim\dzg@ias@ell{\l}pt>0pt
          \pgfmathparse{abs(\csname dzg@ias@ax@\apex\endcsname
              -\csname dzg@ias@ax@\bapex\endcsname)
            +abs(\csname dzg@ias@ay@\apex\endcsname
              -\csname dzg@ias@ay@\bapex\endcsname)<0.01 ? 1 : 0}%
          \ifdim\pgfmathresult pt>0pt
            \pgfmathtruncatemacro\dzg@m{\dzg@mult+1}\xdef\dzg@mult{\dzg@m}%
            \ifnum\bapex<\apex \gdef\dzg@first{0}\fi
          \fi
        \fi}%
      \ifnum\dzg@first=1 \iassingularity{-p\apex}{\dzg@mult}\fi
    \fi}%
  % Walk the vertices cyclically from one whose preceding side has positive
  % length, so that a group of collapsed vertices is never split.
  \gdef\dzg@start{#1}%
  \ifnum#3=1
    \gdef\dzg@found{0}%
    \foreach \i in {#1,...,#2} {%
      \pgfmathtruncatemacro\dzg@prev{\i==#1 ? #2 : \i-1}%
      \ifnum\dzg@found=0 \ifdim\csname dzg@ias@len@\dzg@prev\endcsname pt>0pt
        \xdef\dzg@start{\i}\gdef\dzg@found{1}\fi\fi}%
  \fi
  \gdef\dzg@mult{0}%
  \pgfmathtruncatemacro\dzg@last{\dzg@ias@n-1}%
  \foreach \k in {0,...,\dzg@last} {%
    \pgfmathtruncatemacro\i{#1+mod(\dzg@start-#1+\k,\dzg@ias@n)}%
    \gdef\dzg@charge{1}%
    \foreach \e/\m in {#6} {\ifnum\e=\i \xdef\dzg@charge{\m}\fi}%
    \pgfmathtruncatemacro\dzg@m{\dzg@mult+\dzg@charge}\xdef\dzg@mult{\dzg@m}%
    \gdef\dzg@emit{0}%
    \ifnum#3=0 \ifnum\i=#2 \gdef\dzg@emit{1}\fi\fi
    \ifnum\dzg@emit=0 \ifnum\i<#2
      \ifdim\csname dzg@ias@len@\i\endcsname pt>0pt \gdef\dzg@emit{1}\fi\fi\fi
    \ifnum#3=1 \ifnum\i=#2
      \ifdim\csname dzg@ias@len@\i\endcsname pt>0pt \gdef\dzg@emit{1}\fi\fi\fi
    \ifnum\dzg@emit=1
      \iassingularity{-p\i}{\dzg@mult}%
      \ifiaslabels
        \pgfmathsetmacro\dzg@a{atan2(\csname dzg@ias@y@\i\endcsname-\dzg@ias@cy,
          \csname dzg@ias@x@\i\endcsname-\dzg@ias@cx)}%
        % Inside the polygon, so that point names and edge vectors do not meet.
        \path (-p\i) node[ias label, anchor=\dzg@a, inner sep=6pt] {$p_{\i}$};
      \fi
      \gdef\dzg@mult{0}%
    \fi}}


\pgfkeys{/dzg/symington/.cd, first/.store in=\dzg@sym@first,
  last/.store in=\dzg@sym@last, closed/.code={\def\dzg@sym@closed{1}},
  charges/.store in=\dzg@sym@charges}
\newcommand\symingtonpolygon[3][]{%
  \def\dzg@sym@first{0}\def\dzg@sym@last{0}\def\dzg@sym@closed{0}%
  \def\dzg@sym@charges{-1/1}%
  \pgfkeys{/dzg/symington/.cd, #1}%
  \def\dzg@sym@sides{#2}\def\dzg@sym@surgeries{#3}%
  \edef\dzg@next{\noexpand\dzg@ias@symington{\dzg@sym@first}{\dzg@sym@last}%
    {\dzg@sym@closed}\noexpand\dzg@sym@sides\noexpand\dzg@sym@surgeries
    {\dzg@sym@charges}}%
  \dzg@next}

% \setell{<first index>}{<default values>}: l_i from the figure's `ell`, else
% the object's default instance. \ellvalue{<i>} is l_i.
\newcommand\setell[2]{%
  \ifx\dzg@ell@given\relax\dzg@ias@setell{#1}{#2}%
  \else\edef\dzg@next{\noexpand\dzg@ias@setell{#1}{\dzg@ell@given}}\dzg@next\fi}
\newcommand\ellvalue[1]{\dzg@ias@ell{#1}}

% =============================================================================
% ADE lattice polygons
% =============================================================================
% \adepolygon[ade affine]{<px>/<py>}{<x/y, ...>}: the lattice polygon Q with
% boundary p* -> v_1 -> ... -> v_k -> p* counterclockwise, so that
% C = [v_k, p*] + [p*, v_1] and C' = v_1 -> ... -> v_k (Alexeev-Thompson,
% "ADE surfaces and their moduli", arXiv:1712.07932, Notation 3.7). With
% `ade affine` the boundary is v_1 -> ... -> v_k -> v_1 and p* is interior to
% the side [v_k, v_1], which is C. What it draws, all derived from Q and p*:
% - the lattice [0,w]x[0,h] with points named -<i>-<j>, and Q filled;
% - C as `long side` (even lattice length) or `short side` (odd), labelled
%   with its lattice length; the other sides C' as `axis side`; the lattice
%   points of C as `boundary point` in the accent colour, and p* (named
%   -pstar) as `marked point`;
% - the Dynkin diagram of AT Notation 3.7: a node at each lattice point of C'
%   not on C, joined along the boundary, black at lattice length 2 from p* and
%   white at length 1; at each corner node an internal node (`vertex=doubled`)
%   at the corner plus the primitive vectors of its two sides. Nodes are named
%   -1, -2, ... along the boundary, then the internal nodes;
% - the monomial x^i y^j at every boundary lattice point but p*, and at every
%   internal node, set outside Q (along the outward normal of the side, or the
%   bisector of the two normals at a vertex) and, at an internal node, away
%   from its corner.
% The local bounding box of Q is -lp.

% \dzg@ade@side{<x0>}{<y0>}{<x1>}{<y1>}: lattice length \dzg@g and primitive
% step (\dzg@sx,\dzg@sy) of the segment.
\def\dzg@ade@side#1#2#3#4{%
  \pgfmathtruncatemacro\dzg@g{gcd(abs(#3-#1),abs(#4-#2))}%
  \pgfmathtruncatemacro\dzg@sx{(#3-#1)/\dzg@g}%
  \pgfmathtruncatemacro\dzg@sy{(#4-#2)/\dzg@g}}

% \dzg@ade@monomial{<i>}{<j>}{<angle>}
\def\dzg@ade@monomial#1#2#3{%
  \ifadelabels
    \path (#1,#2) node[label={[monomial]#3:$\dzgmonomial{#1}{#2}$}] {};%
  \fi}

\newcommand\adepolygon[3][]{%
  \begingroup\tikzset{#1}\dzg@adepolygon#2\relax{#3}\endgroup}
\def\dzg@adepolygon#1/#2\relax#3{%
  \def\dzg@px{#1}\def\dzg@py{#2}\def\dzg@path{#3}%
  \ifadeaffine \let\dzg@cycle\dzg@path
  \else\edef\dzg@cycle{\dzg@px/\dzg@py, \dzg@path}\fi
  % Bounding box, the path of Q, and v_1, v_k.
  \gdef\dzg@w{0}\gdef\dzg@h{0}\gdef\dzg@Q{}%
  \foreach \x/\y [count=\k] in \dzg@cycle {%
    \xdef\dzg@N{\k}%
    \pgfmathtruncatemacro\dzg@t{max(\dzg@w,\x)}\xdef\dzg@w{\dzg@t}%
    \pgfmathtruncatemacro\dzg@t{max(\dzg@h,\y)}\xdef\dzg@h{\dzg@t}%
    \xdef\dzg@Q{\dzg@Q(\x,\y) -- }}%
  \foreach \x/\y [count=\k] in \dzg@path {%
    \ifnum\k=1 \xdef\dzg@vfirst{\x/\y}\xdef\dzg@vfirstx{\x}\xdef\dzg@vfirsty{\y}\fi
    \xdef\dzg@vlast{\x/\y}\xdef\dzg@vlastx{\x}\xdef\dzg@vlasty{\y}}%
  \path[local bounding box=-lp] (0,0) rectangle (\dzg@w,\dzg@h);
  \begin{scope}[on background layer]
    \fill[polygon region] \dzg@Q cycle;
    \draw[lattice grid] (0,0) grid (\dzg@w,\dzg@h);
  \end{scope}
  \foreach \x in {0,...,\dzg@w} \foreach \y in {0,...,\dzg@h}
    \node[lattice point] (-\x-\y) at (\x,\y) {};
  % Monomials: outward along the normal of each side at its interior points,
  % along the bisector of the two outward normals at each vertex.
  \edef\dzg@wrapped{\dzg@cycle, \dzg@cycle}%
  \foreach \x/\y [count=\k, remember=\xa as \xb (initially 0),
      remember=\ya as \yb (initially 0), remember=\x as \xa (initially 0),
      remember=\y as \ya (initially 0)] in \dzg@wrapped {%
    \pgfmathtruncatemacro\dzg@atvertex{\k>2 && \k<\dzg@N+3}%
    \pgfmathtruncatemacro\dzg@atside{\k>1 && \k<\dzg@N+2}%
    \ifnum\dzg@atvertex=1 % vertex (\xa,\ya) between (\xb,\yb) and (\x,\y)
      \pgfmathsetmacro\dzg@nx{(\ya-\yb)/veclen(\xa-\xb,\ya-\yb)
        + (\y-\ya)/veclen(\x-\xa,\y-\ya)}%
      \pgfmathsetmacro\dzg@ny{(\xb-\xa)/veclen(\xa-\xb,\ya-\yb)
        + (\xa-\x)/veclen(\x-\xa,\y-\ya)}%
      \pgfmathsetmacro\dzg@angle{atan2(\dzg@ny,\dzg@nx)}%
      \ifnum\xa=\dzg@px \ifnum\ya=\dzg@py \else
        \dzg@ade@monomial{\xa}{\ya}{\dzg@angle}\fi
      \else\dzg@ade@monomial{\xa}{\ya}{\dzg@angle}\fi
    \fi
    \ifnum\dzg@atside=1 % interior lattice points of side (\xa,\ya) -> (\x,\y)
      \dzg@ade@side{\xa}{\ya}{\x}{\y}%
      \pgfmathsetmacro\dzg@angle{atan2(-\dzg@sx,\dzg@sy)}%
      \pgfmathtruncatemacro\dzg@last{\dzg@g-1}%
      \ifnum\dzg@last>0
        \foreach \t in {1,...,\dzg@last} {%
          \pgfmathtruncatemacro\dzg@i{\xa+\t*\dzg@sx}%
          \pgfmathtruncatemacro\dzg@j{\ya+\t*\dzg@sy}%
          \ifnum\dzg@i=\dzg@px \ifnum\dzg@j=\dzg@py \else
            \dzg@ade@monomial{\dzg@i}{\dzg@j}{\dzg@angle}\fi
          \else\dzg@ade@monomial{\dzg@i}{\dzg@j}{\dzg@angle}\fi}%
      \fi
    \fi}%
  % C: the sides through p*, and their lattice points.
  \ifadeaffine
    \edef\dzg@Csides{\dzg@vlast/\dzg@vfirst}%
  \else
    \edef\dzg@Csides{\dzg@vlast/\dzg@px/\dzg@py, \dzg@px/\dzg@py/\dzg@vfirst}%
  \fi
  \foreach \xa/\ya/\x/\y in \dzg@Csides {%
    \dzg@ade@side{\xa}{\ya}{\x}{\y}%
    \ifodd\dzg@g\def\dzg@class{short side}\else\def\dzg@class{long side}\fi
    \scoped[on edge layer] \draw[\dzg@class] (\xa,\ya) -- (\x,\y);
    \ifadelabels \path[side label=\dzg@g] (\xa,\ya) -- (\x,\y);\fi
    \foreach \t in {0,...,\dzg@g}
      \node[boundary point=dzg accent] at (\xa+\t*\dzg@sx,\ya+\t*\dzg@sy) {};}%
  \coordinate (-pstar) at (\dzg@px,\dzg@py);
  \ifadelabels \node[marked point=$p_\ast$] at (-pstar) {};
  \else \node[marked point] at (-pstar) {};\fi
  % C' and the Dynkin diagram along it.
  \gdef\dzg@node{0}%
  \foreach \x/\y [count=\k, remember=\x as \xa (initially 0),
      remember=\y as \ya (initially 0)] in \dzg@path {%
    \ifnum\k>1
      \scoped[on edge layer] \draw[axis side] (\xa,\ya) -- (\x,\y);
      \dzg@ade@side{\xa}{\ya}{\x}{\y}%
      \foreach \t in {1,...,\dzg@g} {%
        \pgfmathtruncatemacro\dzg@i{\xa+\t*\dzg@sx}%
        \pgfmathtruncatemacro\dzg@j{\ya+\t*\dzg@sy}%
        \edef\dzg@here{\dzg@i/\dzg@j}%
        \ifx\dzg@here\dzg@vlast\else
          \pgfmathtruncatemacro\dzg@n{\dzg@node+1}\xdef\dzg@node{\dzg@n}%
          \pgfmathtruncatemacro\dzg@len{gcd(abs(\dzg@i-\dzg@px),abs(\dzg@j-\dzg@py))}%
          \ifnum\dzg@len=2 \def\dzg@mark{black}\else\def\dzg@mark{white}\fi
          \node[vertex=\dzg@mark] (-\dzg@n) at (\dzg@i,\dzg@j) {};
          \ifnum\dzg@n>1
            \scoped[on edge layer]
              \draw[coxeter edge=3] (\dzg@i-\dzg@sx,\dzg@j-\dzg@sy) -- (\dzg@i,\dzg@j);
          \fi
        \fi}%
    \fi}%
  % Internal nodes at the corners of C'.
  \foreach \x/\y [count=\k, remember=\xa as \xb (initially 0),
      remember=\ya as \yb (initially 0), remember=\x as \xa (initially 0),
      remember=\y as \ya (initially 0)] in \dzg@path {%
    \ifnum\k>2 % corner (\xa,\ya)
      \dzg@ade@side{\xa}{\ya}{\xb}{\yb}\let\dzg@ux\dzg@sx\let\dzg@uy\dzg@sy
      \dzg@ade@side{\xa}{\ya}{\x}{\y}%
      \pgfmathtruncatemacro\dzg@dx{\dzg@ux+\dzg@sx}%
      \pgfmathtruncatemacro\dzg@dy{\dzg@uy+\dzg@sy}%
      \pgfmathtruncatemacro\dzg@i{\xa+\dzg@dx}%
      \pgfmathtruncatemacro\dzg@j{\ya+\dzg@dy}%
      \pgfmathtruncatemacro\dzg@n{\dzg@node+1}\xdef\dzg@node{\dzg@n}%
      \node[vertex=doubled] (-\dzg@n) at (\dzg@i,\dzg@j) {};
      \scoped[on edge layer] \draw[coxeter edge=3] (\xa,\ya) -- (\dzg@i,\dzg@j);
      \pgfmathsetmacro\dzg@angle{atan2(\dzg@dy,\dzg@dx)}%
      \dzg@ade@monomial{\dzg@i}{\dzg@j}{\dzg@angle}%
    \fi}}

% \adetitle{<math>}: the decorated Dynkin symbol (AT Notations 3.8, 3.21) under
% the polygon drawn last.
\newcommand\adetitle[1]{%
  \ifadetitle\node[polygon title, below=7mm of -lp.south] {#1};\fi}

% =============================================================================
% Complete fans of toric varieties
% =============================================================================
% \fanrays{<i/a/b/label anchor, ...>}: the rays from the origin -0 to -rho<i>,
% the primitive generator (a,b) times \fanlength, labelled rho_i = (a,b).
% \fancones{<i/coordinate, ...>}: the maximal cones sigma_i, labelled at
% -sigma<i>. Both respect `fan labels`.
\newcommand\fanrays[1]{%
  \node[distinguished point] (-0) at (0,0) {};
  \foreach \dzg@i/\dzg@a/\dzg@b/\dzg@p in {#1}{%
    \coordinate (-rho\dzg@i) at (\fanlength*\dzg@a, \fanlength*\dzg@b);
    \iffanlabels
      \draw[fan ray] (-0) -- (-rho\dzg@i)
        node[\dzg@p, font=\small] {$\rho_{\dzg@i} = (\dzg@a,\dzg@b)$};
    \else
      \draw[fan ray] (-0) -- (-rho\dzg@i);
    \fi}}
\newcommand\fancones[1]{%
  \foreach \dzg@i/\dzg@c in {#1}{%
    \coordinate (-sigma\dzg@i) at (\dzg@c);
    \iffanlabels \node[font=\small] at (-sigma\dzg@i) {$\sigma_{\dzg@i}$};\fi}}

% =============================================================================
% Coxeter diagrams: constructors
% =============================================================================
% Building blocks that draw a Coxeter diagram from its data: roots (vertices)
% with their marks and positions, edges with their Coxeter weights, chains,
% cycles and maximal parabolic subdiagrams. The diagrams themselves are
% objects, one file each, under figures/objects/coxeter/, assembled from these
% constructors; a figure places one with
%   \pic (A) at (6,0) [<keys>] {object=coxeter/sterk-cusp-1};
% and draws on its roots A-0, A-1, ...:
%   \scoped[on background layer] \draw[parabolic] (A-0.center) -- (A-8.center);
% \input by dzg-constructors.tex, so "@" is a letter here.
%
% Commands (the root <i> is the node -<i>, labelled by the label scheme):
%   \coxeterlabels{<scheme>}
%       the object's canonical label scheme, used unless the figure passes
%       `root labels`; an object states it before its first root.
%   \coxeterroot{<i>}{<mark>}{<coordinate>}{<label direction>}
%       one root. <mark> is a `vertex` style (white, black, doubled) or even.
%   \coxeterroots{<i>/<mark>/<coordinate>/<label direction>, ...}
%       several roots. Brace a coordinate that contains a comma; place a calc
%       coordinate with \coxeterroot, since \foreach keeps the braces.
%   \coxeteredges{<Coxeter weight>}{<path through -<i> names>}
%       edges of one weight (a `coxeter edge` value), on the edge layer.
%   \coxeterchain{<Coxeter weight>}{<first>}{<last>}
%       the edges i -- i+1 for first <= i < last.
%   \coxetersquare{<side>}{<edges per side>}{<first>}{<sw|nw>}
%       {<mark at corners and every second root>}{<mark between>}
%       a cycle of simple edges around the square [0,side]^2, numbered from
%       <first> at the south-west corner counterclockwise (sw) or at the
%       north-west corner clockwise (nw), labelled outward.
%   \coxeterEdiagram{<n>}{<mark>}   (dzg-vinberg.tex)
%   \coxeterparabolic{<name>}{{<i>,<j>,...}, {...}, ...}
%       a maximal parabolic subdiagram of the object: the highlight stroked
%       through the root centres of each chain, drawn when the figure passes
%       `maximal parabolic=<name>`. The vertex sets are data of the object.
%   \coxeterpair[<keys>], \coxetervertex{<mark>}   table legends
%
% Pics:
%   dynkin folding={<unfolded options>}{<diagram>}{<folded options>}{<diagram>}
%   root chain={<mark>/<norm label>/<weight to the next root>, ...}
%
% Keys:
%   root labels=alpha|r|ell|index|none   the label scheme (see below); a bare
%                                        `root labels` selects the canonical one
%   parity marks=true|false              even roots doubled, or white
%   maximal parabolic=<name>             highlight that parabolic subdiagram
%   vertex=even                          an even boundary root (AEGS25
%                                        Sec. 3.5): doubled under `parity
%                                        marks`, white otherwise
%
% Label schemes; the index is always the canonical root name:
%   alpha  $\alpha_i$
%   r      $r_i$
%   ell    $\ell_i$    the coordinate \ell_i = \lambda\cdot\alpha_i of an
%                      integral-affine sphere B(\lambda) (IAS figures)
%   index  $i$         the bare index
%   none   no labels

\newif\ifdzg@rootlabels
\newif\ifdzg@paritymarks
\dzg@paritymarkstrue
\def\dzg@unset{unset}
\def\dzg@labelscheme{unset}
\def\dzg@parabolic{}
\tikzset{
  root labels/.is choice,
  root labels/canonical/.code={\dzg@rootlabelstrue\def\dzg@labelscheme{unset}},
  root labels/alpha/.code={\dzg@rootlabelstrue\def\dzg@labelscheme{alpha}},
  root labels/r/.code={\dzg@rootlabelstrue\def\dzg@labelscheme{r}},
  root labels/ell/.code={\dzg@rootlabelstrue\def\dzg@labelscheme{ell}},
  root labels/index/.code={\dzg@rootlabelstrue\def\dzg@labelscheme{index}},
  root labels/none/.code={\dzg@rootlabelsfalse\def\dzg@labelscheme{none}},
  root labels/.default=canonical,
  parity marks/.is if=dzg@paritymarks,
  vertex/even/.code={%
    \ifdzg@paritymarks\pgfkeysalso{vertex/doubled}\else\pgfkeysalso{vertex/white}\fi},
  maximal parabolic/.code={\def\dzg@parabolic{#1}},
}
\def\dzg@label@alpha#1{$\alpha_{#1}$}
\def\dzg@label@r#1{$r_{#1}$}
\def\dzg@label@ell#1{$\ell_{#1}$}
\def\dzg@label@index#1{$#1$}
\def\dzg@label@unset#1{$r_{#1}$}
\def\dzg@rootlabel#1{\csname dzg@label@\dzg@labelscheme\endcsname{#1}}

\newcommand\coxeterlabels[1]{%
  \ifx\dzg@labelscheme\dzg@unset\dzg@rootlabelstrue\def\dzg@labelscheme{#1}\fi}

\newcommand\coxeterroot[4]{%
  \ifdzg@rootlabels
    \node[vertex=#2, label={[vertex label]#4:\dzg@rootlabel{#1}}] (-#1) at (#3) {};%
  \else
    \node[vertex=#2] (-#1) at (#3) {};%
  \fi}

\newcommand\coxeterroots[1]{%
  \foreach \dzg@i/\dzg@m/\dzg@p/\dzg@a in {#1}{\coxeterroot{\dzg@i}{\dzg@m}{\dzg@p}{\dzg@a}}}

\newcommand\coxeteredges[2]{\scoped[on edge layer] \draw[coxeter edge=#1] #2;}

% A path \foreach through node names restarts every segment at the first node,
% so each edge is its own path.
\newcommand\coxeterchain[3]{%
  \pgfmathtruncatemacro\dzg@stop{#3-1}%
  \foreach \dzg@k in {#2,...,\dzg@stop}{%
    \pgfmathtruncatemacro\dzg@n{\dzg@k+1}%
    \coxeteredges{#1}{(-\dzg@k) -- (-\dzg@n)}}}

\newcommand\coxetersquare[6]{%
  \def\dzg@L{#1}%
  \pgfmathtruncatemacro\dzg@last{#3+4*#2-1}%
  \foreach \dzg@k in {#3,...,\dzg@last}{%
    \pgfmathtruncatemacro\dzg@s{div(\dzg@k-#3,#2)}%
    \pgfmathtruncatemacro\dzg@t{mod(\dzg@k-#3,#2)}%
    \pgfmathsetmacro\dzg@u{#1*\dzg@t/#2}%
    \pgfmathsetmacro\dzg@v{#1-\dzg@u}%
    \csname dzg@square@#4\endcsname
    \ifodd\dzg@t\def\dzg@m{#6}\else\def\dzg@m{#5}\fi
    \coxeterroot{\dzg@k}{\dzg@m}{\dzg@p}{\dzg@a}}%
  \coxeterchain{3}{#3}{\dzg@last}%
  \coxeteredges{3}{(-\dzg@last) -- (-#3)}}
% Position \dzg@p and outward label direction \dzg@a of the root at step
% \dzg@t along side \dzg@s.
\def\dzg@square@corner#1#2{\ifnum\dzg@t=0 \def\dzg@a{#1}\else\def\dzg@a{#2}\fi}
\def\dzg@square@sw{%
  \ifcase\dzg@s
    \edef\dzg@p{\dzg@u,0}\dzg@square@corner{225}{270}%
  \or\edef\dzg@p{\dzg@L,\dzg@u}\dzg@square@corner{315}{0}%
  \or\edef\dzg@p{\dzg@v,\dzg@L}\dzg@square@corner{45}{90}%
  \or\edef\dzg@p{0,\dzg@v}\dzg@square@corner{135}{180}%
  \fi}
\def\dzg@square@nw{%
  \ifcase\dzg@s
    \edef\dzg@p{\dzg@u,\dzg@L}\dzg@square@corner{135}{90}%
  \or\edef\dzg@p{\dzg@L,\dzg@v}\dzg@square@corner{45}{0}%
  \or\edef\dzg@p{\dzg@v,0}\dzg@square@corner{315}{270}%
  \or\edef\dzg@p{0,\dzg@u}\dzg@square@corner{225}{180}%
  \fi}

% Each chain is its own stroke, starting with a zero-length segment so that a
% chain of one root is a dot.
\newcommand\coxeterparabolic[2]{%
  \def\dzg@name{#1}%
  \ifx\dzg@name\dzg@parabolic
    \scoped[on background layer]{%
      \foreach \dzg@c in {#2} {%
        \foreach \dzg@v [count=\dzg@k] in \dzg@c {%
          \ifnum\dzg@k=1 \xdef\dzg@path{(-\dzg@v.center) -- (-\dzg@v.center)}%
          \else\xdef\dzg@path{\dzg@path -- (-\dzg@v.center)}\fi}%
        \draw[parabolic] \dzg@path;}}%
  \fi}

% -----------------------------------------------------------------------------
% Folding of Dynkin diagrams (dynkin-diagrams package): the diagram with its
% symmetry above the folded diagram, joined by the quotient map.
%   \pic {dynkin folding={<unfolded options>}{<type><spec>}{<folded options>}
%     {<type><spec>}};
% e.g. {ply=2}{A{ooo.ooooooo.ooo}}{}{C{ooo.ooo*}}. The nodes -unfolded and
% -folded are the two diagrams.
% -----------------------------------------------------------------------------
% The folded diagram is left-aligned under the unfolded one. Its node resets
% `anchor`, which the nested picture of \dynkin would otherwise inherit.
\tikzset{pics/dynkin folding/.style n args={4}{code={%
  \node[inner sep=2pt] (-unfolded) at (0,0) {\dynkin[scale=3,#1]#2};
  \node[inner sep=2pt, anchor=north west] (-folded)
    at ($(-unfolded.south west)+(0,-10mm)$)
    {\tikzset{anchor=center}\dynkin[scale=3, arrows=false, #3]#4};
  \draw[quotient map] (-unfolded.south) -- (-unfolded.south |- -folded.north);}}}

% -----------------------------------------------------------------------------
% A linear Coxeter diagram with the norm of each root written above it:
%   \pic {root chain={<mark>/<norm label>/<weight to the next root>, ...,
%     <mark>/<norm label>/}};
% Weight 3 is a simple edge; the vertices are named -0, -1, ....
% -----------------------------------------------------------------------------
\tikzset{pics/root chain/.style={code={%
  \foreach \dzg@m/\dzg@l/\dzg@w [count=\dzg@k from 0] in {#1}{%
    \node[vertex=\dzg@m, label={[vertex label]90:{$\dzg@l$}}] (-\dzg@k) at (1.5*\dzg@k,0) {};
    \ifx\dzg@w\pgfutil@empty\else
      \coordinate (-\dzg@k-next) at (1.5*\dzg@k+1.5,0);
      \scoped[on edge layer] \draw[coxeter edge=\dzg@w] (-\dzg@k.center) -- (-\dzg@k-next);
    \fi}}}}

% -----------------------------------------------------------------------------
% Table legends: two roots and the edge between them, or one vertex.
% -----------------------------------------------------------------------------
%   \coxeterpair[marks=black/white, edge=4, roots=f_1/f_2, weight=m]
% marks: the two vertex marks (default white/white); edge: a coxeter edge
% value or none (default 3); roots: the two labels in math mode (default
% H_1/H_2); weight: text set above the edge (default empty).
%   \coxetervertex{white|black|doubled}
\pgfqkeys{/dzg/coxeter pair}{
  marks/.store in=\dzg@pairmarks, marks=white/white,
  edge/.store in=\dzg@pairedge, edge=3,
  roots/.store in=\dzg@pairroots, roots=H_1/H_2,
  weight/.store in=\dzg@pairweight, weight=,
}
\def\dzg@none{none}
\def\dzg@splitpair#1/#2\@nil#3#4{\def#3{#1}\def#4{#2}}
\newcommand{\coxeterpair}[1][]{%
  \begin{tikzpicture}[baseline=-0.5ex]
    \pgfqkeys{/dzg/coxeter pair}{#1}%
    \expandafter\dzg@splitpair\dzg@pairmarks\@nil\dzg@marka\dzg@markb
    \expandafter\dzg@splitpair\dzg@pairroots\@nil\dzg@roota\dzg@rootb
    \node[vertex=\dzg@marka, label={[vertex label]above:$\dzg@roota$}] (h1) at (0,0) {};
    \node[vertex=\dzg@markb, label={[vertex label]above:$\dzg@rootb$}] (h2) at (3,0) {};
    \ifx\dzg@pairedge\dzg@none\else
      \scoped[on edge layer] \draw[coxeter edge=\dzg@pairedge]
        (h1) -- node[vertex label, midway, above] {\dzg@pairweight} (h2);
    \fi
  \end{tikzpicture}}
\newcommand{\coxetervertex}[1]{\tikz[baseline=-0.5ex]\node[vertex=#1] {};}

% =============================================================================
% Vinberg diagrams and mirror moves: constructors
% =============================================================================
% Built on the root, edge and label constructors of dzg-coxeter.tex. The
% objects are figures under figures/objects/coxeter/ (Coxeter and Vinberg
% diagrams) and figures/objects/mirror-moves/ (mirror-move posets).
% \input by dzg-constructors.tex, so "@" is a letter here.
%
% Commands:
%   \coxeterEdiagram{<n>}{<mark>}
%       the tree T_{2,3,n-3} with every root marked <mark>, in Bourbaki's
%       numbering of E_n: the chain 1,3,4,...,n on the x-axis (root k >= 3
%       at x = k-2, root 1 at 0), root 2 above root 4, labels below.
%   \latticevertex[<node options>][<label options>]{<name>}{<mark>}
%       {<coordinate>}{<lattice>}{<label direction>}
%       a lattice in a mirror-move poset: the node -<name>, a vertex marked as
%       in Nikulin's table (white for delta = 0, black for delta = 1),
%       labelled $<lattice>$.
%   \mirrorstages{<y>}
%       the column heads S, T = S^perp, \bar T, \bar{\bar T} at x = 0, 2, 4, 6,
%       nodes -stage0, -stage2, -stage4, -stage6.
% An arrow is a move, styled by its type: `mirror move=complement|U|U(2)`.
%
% Keys:
%   mirror targets=valid|all   draw only the moves to lattices of valid cusps
%                              (valid) or every target of every move (all).
% Read by objects:
%   \ifmirroralltargets        true under `mirror targets=all`
%   \mirrorvalidcolor          text colour of a valid lattice: dzg accent
%                              under `mirror targets=all`, the current colour
%                              otherwise
%   \mirrorvalidlabel          label direction of a valid lattice: below left
%                              under `mirror targets=all`, below otherwise

\newcommand\coxeterEdiagram[2]{%
  \coxeterroot{1}{#2}{0,0}{below}%
  \coxeterroot{2}{#2}{2,1}{above}%
  \foreach \dzg@k in {3,...,#1}{%
    \pgfmathtruncatemacro\dzg@x{\dzg@k-2}%
    \coxeterroot{\dzg@k}{#2}{\dzg@x,0}{below}}%
  \coxeteredges{3}{(-1) -- (-3) (-2) -- (-4)}%
  \coxeterchain{3}{3}{#1}}

\NewDocumentCommand\latticevertex{O{} O{} m m m m m}{%
  \node[#1, vertex=#4, label={[vertex label, #2]#7:$#6$}] (-#3) at (#5) {};}

\newcommand\mirrorstages[1]{%
  \foreach \dzg@x/\dzg@s in {0/S, 2/T, 4/\bar T, 6/\bar{\bar T}}
    \node (-stage\dzg@x) at (\dzg@x, #1) {$\dzg@s$};}

\newif\ifmirroralltargets
\tikzset{
  mirror targets/.is choice,
  mirror targets/valid/.code={\mirroralltargetsfalse
    \def\mirrorvalidcolor{.}\def\mirrorvalidlabel{below}},
  mirror targets/all/.code={\mirroralltargetstrue
    \def\mirrorvalidcolor{dzg accent}\def\mirrorvalidlabel{below left}},
  mirror targets=valid,
}

% =============================================================================
% Baily-Borel cusp diagrams: constructors
% =============================================================================
% Building blocks for cusp diagrams of Baily-Borel compactifications. A
% constructor takes the data of cusps and incidences and draws them; the cusp
% diagrams themselves are objects, one file each, under figures/objects/bb-cusps/,
% assembled from these constructors. A 0-cusp (the orbit of a primitive
% isotropic line) is a cusp0 node and a 1-cusp (the orbit of a primitive
% isotropic plane) a cusp1 node; an incidence arrow runs from a 0-cusp to each
% 1-cusp whose closure contains it. A 0-cusp or 1-cusp is named by its boundary
% lattice I^perp/I; where the lattice has a standard name it is used (E_{10}(2)
% = U(2)+E_8(2), D_8^*(2), D_{16}^+, A_1^{+7}). Place a diagram with
%   \pic (S) {object=bb-cusps/fen2};
% and decorate through the node centres, e.g. the image of another boundary:
%   \scoped[on background layer] \draw[cusp image] (S-2.center) -- (S-245.center);
% \input by dzg-constructors.tex, so "@" is a letter here.
%
% Commands:
%   \bbcusp{<cusp0|cusp1, options>}{<name>}{<coordinate>}{<content>}
%          {<label direction>}{<canonical label>}{<eta label>}
%       a cusp node with content $<content>$ and the label of the active
%       scheme, $<canonical label>$ or $<eta label>$; an empty label in the
%       active scheme draws none.
%   \bblabel{<canonical>}{<eta>}   the text of the active label scheme (empty
%                                  under cusp labels=none), for node contents
%   \bbincidences{<0-cusp>/<1-cusp>/<divisibility>, ...}
%       incidence arrows; divisibility is that of the second isotropic vector
%       of the plane in the 0-cusp lattice, 1 or 2 (an even incidence).
%   \bbbrace{<x from>}{<x to>}{<y>}{<text>}
%       a brace under a column with the invariants $<text>$ of its lattices,
%       drawn in the canonical scheme only.
%   \bbcolumn{<x>}{<name prefix>}{<first>}{<second>}{<next to last>}{<last>}
%       a column of four 1-cusps at x, y = 2, 1, -1, -2, named <prefix>1, ...,
%       <prefix>4, with a vertical ellipsis at y = 0: a family C_1, ..., C_n.
%
% Keys:
%   cusp labels=canonical   the canonical labels of each object (default)
%   cusp labels=eta         the names eta_i, I_{i,j} of the dissertation text
%                           (Enriques-K3 chapter) where the object has them
%   cusp labels=none        no labels
%   incidence divisibility  double each incidence whose second isotropic
%                           vector has divisibility 2 in the 0-cusp lattice
%                           (an even move); a single arrow is divisibility 1.
%                           For 2-elementary lattices this is the incidence at
%                           which the a-invariant drops by 2 [AE22, Prop. 5.5].

\def\dzg@bb@scheme{canonical}
\tikzset{
  cusp labels/.is choice,
  cusp labels/canonical/.code={\def\dzg@bb@scheme{canonical}},
  cusp labels/eta/.code={\def\dzg@bb@scheme{eta}},
  cusp labels/none/.code={\def\dzg@bb@scheme{none}},
  incidence divisibility/.style={even incidence/.style={doubled mark}},
  even incidence/.style={},
}

\def\dzg@bb@canonical{canonical}
\def\dzg@bb@eta{eta}
\def\dzg@bb@none{none}
\newcommand\bblabel[2]{\csname dzg@bb@label@\dzg@bb@scheme\endcsname{#1}{#2}}
\def\dzg@bb@label@canonical#1#2{#1}
\def\dzg@bb@label@eta#1#2{#2}
\def\dzg@bb@label@none#1#2{}

\newcommand\bbcusp[7]{%
  \ifx\dzg@bb@scheme\dzg@bb@canonical\def\dzg@bb@text{#6}\fi
  \ifx\dzg@bb@scheme\dzg@bb@eta\def\dzg@bb@text{#7}\fi
  \ifx\dzg@bb@scheme\dzg@bb@none\let\dzg@bb@text\@empty\fi
  \ifx\dzg@bb@text\@empty
    \node[#1] (#2) at (#3) {$#4$};%
  \else
    \node[#1, label={[cusp label]#5:$\dzg@bb@text$}] (#2) at (#3) {$#4$};%
  \fi}

\newcommand\bbincidences[1]{%
  \foreach \dzg@p/\dzg@c/\dzg@d in {#1}{%
    \ifnum\dzg@d=2
      \draw[incidence, even incidence] (\dzg@p) -- (\dzg@c);%
    \else
      \draw[incidence] (\dzg@p) -- (\dzg@c);%
    \fi}}

\newcommand\bbbrace[4]{%
  \ifx\dzg@bb@scheme\dzg@bb@canonical
    \draw[cusp brace] (#1,#3) -- (#2,#3) node[midway, below=5.5ex] {$#4$};%
  \fi}

\newcommand\bbcolumn[6]{%
  \foreach \dzg@k/\dzg@y/\dzg@L in {1/2/{#3}, 2/1/{#4}, 3/-1/{#5}, 4/-2/{#6}}
    {\bbcusp{cusp1}{#2\dzg@k}{#1,\dzg@y}{\dzg@L}{}{}{}}%
  \node at (#1,0) {$\vdots$};}

% =============================================================================
% Tables of lattice invariants: constructor
% =============================================================================
% A table of lattice invariants is the grid of integer points (x, y) of two
% invariants, with a vertex mark at each point where a lattice with those
% invariants occurs; the mark says which values of a third invariant occur
% there. Nikulin's table of 2-elementary lattices by (r, a, delta),
% figures/objects/nikulin/two-elementary-table.tikz, is one such object.
% \input by dzg-constructors.tex, so "@" is a letter here.
%
% Command:
%   \latticetable{<x name>}{<x max>}{<y name>}{<y max>}{<x/y/mark, ...>}
%       the grid [0, x max + 1] x [0, y max + 1] with axes named <x name> and
%       <y name> and ticks 1, ..., x max and 1, ..., y max; at each listed
%       point a node[vertex=<mark>] (mark white, black or doubled) named
%       -<x>-<y>.
%
% Key:
%   lattice table labels=axes|none   axis names and ticks (default axes)

\newif\ifdzg@lt@labels
\dzg@lt@labelstrue
\tikzset{
  lattice table labels/.is choice,
  lattice table labels/axes/.code={\dzg@lt@labelstrue},
  lattice table labels/none/.code={\dzg@lt@labelsfalse},
}

\newcommand\latticetable[5]{%
  \pgfmathtruncatemacro\dzg@lt@xgrid{#2+1}%
  \pgfmathtruncatemacro\dzg@lt@ygrid{#4+1}%
  \pgfmathsetmacro\dzg@lt@xend{#2+1.6}%
  \pgfmathsetmacro\dzg@lt@yend{#4+1.6}%
  \draw[lattice grid, step={(1,1)}] (0,0) grid (\dzg@lt@xgrid,\dzg@lt@ygrid);
  \draw[->] (0,0) -- (\dzg@lt@xend,0);
  \draw[->] (0,0) -- (0,\dzg@lt@yend);
  \ifdzg@lt@labels
    \node[right] at (\dzg@lt@xend,0) {$#1$};
    \node[above] at (0,\dzg@lt@yend) {$#3$};
    \foreach \dzg@x in {1,...,#2} {
      \draw (\dzg@x,0.12) -- (\dzg@x,-0.12);
      \node[below=2pt, font=\scriptsize] at (\dzg@x,0) {\dzg@x};
    }
    \foreach \dzg@y in {1,...,#4} {
      \draw (-0.12,\dzg@y) -- (0.12,\dzg@y);
      \node[left=4pt, font=\scriptsize] at (0,\dzg@y) {\dzg@y};
    }
  \fi
  \foreach \dzg@x/\dzg@y/\dzg@m in {#5}
    \node[vertex=\dzg@m] (-\dzg@x-\dzg@y) at (\dzg@x,\dzg@y) {};
}

% =============================================================================
% Moduli-space strata: constructors
% =============================================================================
% Building blocks for the strata of a moduli space of stable curves: dual
% graphs, genus labels, and the poset of strata ordered by closure. The strata
% themselves, their stable graphs and stable curves, are objects, one file
% each, under figures/objects/moduli/, assembled from these constructors; the
% strata of \overline{M}_2 are objects/moduli/m2-*.tikz.
% \input by dzg-constructors.tex, so "@" is a letter here.
%
% Commands:
%   \dualvertex[<label distance>]{<name>}{<x,y>}{<genus>}{<placement>}
%       a vertex of a dual graph with its genus label at <placement>
%       (left, below right, ...) and <label distance> (default 3pt).
%   \dualloop{<x,y>}{left|right}
%       a loop at the vertex at <x,y>: a circle of radius \dualloopradius
%       tangent to the vertex on that side. Draw loops before the vertex so
%       the vertex covers the loop.
%   \genuslabel{<placement>}{<x,y>}{<genus>}
%       the genus of a component of a stable curve, at <x,y>.
%   \strataposet{<object prefix>}{<name/nodes/row, ...>}{<a/b, ...>}
%       the strata of a moduli space as a poset: stratum <name>, with <nodes>
%       nodes (its codimension), at position (nodes, row) of the layout, and
%       a degeneration arrow from a to b for each cover relation a/b
%       (b lies in the closure of a). Stratum <name> is the object
%       <object prefix><render>-<name> placed as -<name>, or with
%       `strata render=name` a poset node labelled <name>; either way its
%       bounding box is -<name>-box, so a stratum object names its bounding
%       box -box. After the call \strataposetdata holds the strata list, for
%       figures that loop over the strata.
%
% Keys:
%   genus labels=all|none       every component and vertex carries its genus,
%                               0 included (default all)
%   strata render=<kind>|name   which object of each stratum \strataposet
%                               places (graph, curve, ...; default graph), or
%                               its name in a poset node
%   strata layout=horizontal|vertical
%                               nodes left to right (x = 4.6 nodes,
%                               y = 1.35 row), or codimension top to bottom
%                               (x = -2.5 row, y = -2 nodes)

\newif\ifdzg@genuslabels
\dzg@genuslabelstrue
\def\dzg@m@render{graph}
\def\dzg@m@layout{horizontal}
\tikzset{
  genus labels/.is choice,
  genus labels/all/.code={\dzg@genuslabelstrue},
  genus labels/none/.code={\dzg@genuslabelsfalse},
  strata render/.store in=\dzg@m@render,
  strata layout/.is choice,
  strata layout/horizontal/.code={\def\dzg@m@layout{horizontal}},
  strata layout/vertical/.code={\def\dzg@m@layout{vertical}},
}

\newcommand\genuslabel[3]{%
  \ifdzg@genuslabels\node[genus, #1] at (#2) {$#3$};\fi}
\newcommand\dualvertex[5][3pt]{%
  \node[dual vertex] (#2) at (#3) {};
  \ifdzg@genuslabels\node[genus, #5=#1] at (#2) {$#4$};\fi}
\def\dualloopradius{0.32}
\def\dzg@m@loop@left{-}
\def\dzg@m@loop@right{}
\newcommand\dualloop[2]{%
  \draw[shift={(#1)}] (\csname dzg@m@loop@#2\endcsname\dualloopradius,0)
    circle (\dualloopradius);}

% Position (\dzg@m@x, \dzg@m@y) of the stratum with <nodes> and <row>.
\def\dzg@m@position#1#2{%
  \csname dzg@m@position@\dzg@m@layout\endcsname{#1}{#2}}
\def\dzg@m@position@horizontal#1#2{%
  \pgfmathsetmacro\dzg@m@x{4.6*#1}\pgfmathsetmacro\dzg@m@y{1.35*#2}}
\def\dzg@m@position@vertical#1#2{%
  \pgfmathsetmacro\dzg@m@x{-2.5*#2}\pgfmathsetmacro\dzg@m@y{-2*#1}}

% \dzg@m@stratum{<object prefix>}{<position>}{<render>}{<name>}
\def\dzg@m@name{name}
\def\dzg@m@stratum#1#2#3#4{%
  \def\dzg@m@r{#3}%
  \ifx\dzg@m@r\dzg@m@name
    \node[poset node] (-#4-box) at (#2) {$#4$};
  \else
    \pic (-#4) at (#2) {object=#1#3-#4};
  \fi}

\newcommand\strataposet[3]{%
  \gdef\strataposetdata{#2}%
  \foreach \dzg@s/\dzg@n/\dzg@row in {#2} {
    \dzg@m@position{\dzg@n}{\dzg@row}
    \edef\dzg@m@next{\noexpand\dzg@m@stratum{#1}{\dzg@m@x,\dzg@m@y}%
      {\dzg@m@render}{\dzg@s}}%
    \dzg@m@next
  }
  \foreach \dzg@a/\dzg@b in {#3}
    \draw[degeneration] (-\dzg@a-box) -- (-\dzg@b-box);
}

% =============================================================================
% Spaces, contours and spectral-sequence pages: constructors
% =============================================================================
% Building blocks for topology and complex-analysis pictures. A constructor
% takes the data of an object and draws it; the objects themselves are files
% under figures/objects/: spaces/ (spheres, balls, cylinder, torus, its
% fundamental square, coffee cup, Moebius band), contours/ (integration
% contours) and diagrams/ (commutative diagrams, a tikz-cd diagram in a node).
% \input by dzg-constructors.tex, so "@" is a letter here.
%
% One decision per former pic or command:
%   sphere=<n>, ball=<n>, cylinder, torus, torus fundamental square, coffee
%     cup, \mobiusband: one particular space each -> objects/spaces/.
%   quarter disc contour, strip rectangle contour: one particular contour
%     each -> objects/contours/.
%   \kleinquotients, \enriquescorrespondence, \cowedgesquare,
%     \monaddescentaxioms: one particular diagram per variant -> objects/diagrams/.
%   space labels, \spacelabel: constructor, the canonical label of any space.
%   \complexaxes: constructor, the axes of C from their extents.
%   ss page, \ssentry, \ssdifferential: constructor, a page from its ranges.
%
% Commands:
%   \spacelabel{<x,y>}{<label>}  the node -label at (x,y), if space labels
%   \complexaxes{<xmin>}{<xmax>}{<ymin>}{<ymax>}
%                                the coordinate axes of C, and the label C at
%                                the top right corner if space labels
%   \pic (E) [<ss keys>] {ss page}, \ssentry, \ssdifferential   see below
%
% Keys:
%   space labels=canonical|none  canonical labels of spaces and contours
%
% Conditional for objects: \ifspacelabels.

\newif\ifspacelabels \spacelabelstrue
\tikzset{
  space labels/.is choice,
  space labels/canonical/.code={\spacelabelstrue},
  space labels/none/.code={\spacelabelsfalse},
}
\newcommand{\spacelabel}[2]{\ifspacelabels\node (-label) at (#1) {#2};\fi}
\newcommand{\complexaxes}[4]{%
  \draw[coordinate axis] (#1,0) -- (#2,0);
  \draw[coordinate axis] (0,#3) -- (0,#4);
  \ifspacelabels\node[font=\small] at (#2-0.3,#4-0.3) {$\mathbb{C}$};\fi}

% -----------------------------------------------------------------------------
% A page of a spectral sequence. The cell (p,q) is the coordinate <name>-p-q at
% (p * ss column sep, q * ss row sep); cells exist for p, q up to three beyond
% the page, so differentials can leave it. The coordinate axes run between
% p = -1 and p = 0 and between q = -1 and q = 0; a page that starts at 0 gets an
% L-shaped frame. Keys (on the pic):
%   ss p={0,...,5}, ss q={0,...,5}   the columns and rows of the page
%   ss p labels={0,n}, ss q labels=  tick labels in place of the indices
%   ss p name=$p$, ss q name=$q$     axis names; ss ticks=false hides ticks
%   ss column sep=2cm, ss row sep=1cm
%   ss homological                   d^r has bidegree (-r, r-1); the default is
%                                    cohomological, d_r of bidegree (r, 1-r)
% Figures place entries and differentials:
%   \ssentry{E}{p}{q}{$E_2^{p,q}$}   the node E-p-q, replacing the coordinate
%   \ssdifferential[<draw options>]{E}{r}{p}{q}
%                                    d_r from (p,q), labelled d_r^{p,q} (or
%                                    d^r_{p,q}); `ss unlabelled` omits the label
% -----------------------------------------------------------------------------
\def\dzg@ssP{0,...,5}\def\dzg@ssQ{0,...,5}
\def\dzg@ssPlabels{}\def\dzg@ssQlabels{}
\def\dzg@ssPname{$p$}\def\dzg@ssQname{$q$}
\def\dzg@ssconv{coh}
\newif\ifdzg@ssticks \dzg@sstickstrue
\newif\ifdzg@sslabel \dzg@sslabeltrue
\tikzset{
  ss p/.store in=\dzg@ssP, ss q/.store in=\dzg@ssQ,
  ss p labels/.store in=\dzg@ssPlabels, ss q labels/.store in=\dzg@ssQlabels,
  ss p name/.store in=\dzg@ssPname, ss q name/.store in=\dzg@ssQname,
  ss column sep/.initial=2cm, ss row sep/.initial=1cm,
  ss ticks/.is if=dzg@ssticks,
  ss homological/.code={\def\dzg@ssconv{hom}},
  ss unlabelled/.code={\dzg@sslabelfalse},
}
% The first and last entries of a foreach list.
\def\dzg@ssends#1#2#3{%
  \def#2{}\foreach \dzg@i in #1 {\ifx#2\empty\xdef#2{\dzg@i}\fi\xdef#3{\dzg@i}}}

\tikzset{pics/ss page/.style={/tikz/transform shape, code={\tikzset{transform shape=false}%
  \pgfmathsetlengthmacro\dzg@sx{\pgfkeysvalueof{/tikz/ss column sep}}%
  \pgfmathsetlengthmacro\dzg@sy{\pgfkeysvalueof{/tikz/ss row sep}}%
  \global\let\dzg@ssx\dzg@sx \global\let\dzg@ssy\dzg@sy
  \dzg@ssends{\dzg@ssP}\dzg@pmin\dzg@pmax
  \dzg@ssends{\dzg@ssQ}\dzg@qmin\dzg@qmax
  \pgfmathtruncatemacro\dzg@pa{\dzg@pmin-3}\pgfmathtruncatemacro\dzg@pb{\dzg@pmax+3}%
  \pgfmathtruncatemacro\dzg@qa{\dzg@qmin-3}\pgfmathtruncatemacro\dzg@qb{\dzg@qmax+3}%
  \foreach \p in {\dzg@pa,...,\dzg@pb} \foreach \q in {\dzg@qa,...,\dzg@qb}
    \coordinate (-\p-\q) at (\p*\dzg@sx, \q*\dzg@sy);
  \expandafter\xdef\csname dzg@ssconv@\pgfkeysvalueof{/tikz/name prefix}\endcsname{\dzg@ssconv}%
  % axes
  \pgfmathsetmacro\dzg@xl{min(\dzg@pmin,0)-0.5}\pgfmathsetmacro\dzg@xr{\dzg@pmax+0.6}%
  \pgfmathsetmacro\dzg@yb{min(\dzg@qmin,0)-0.5}\pgfmathsetmacro\dzg@yt{\dzg@qmax+0.6}%
  \draw (\dzg@xl*\dzg@sx, -0.5*\dzg@sy) -- (\dzg@xr*\dzg@sx, -0.5*\dzg@sy)
    node[right] {\dzg@ssPname};
  \draw (-0.5*\dzg@sx, \dzg@yb*\dzg@sy) -- (-0.5*\dzg@sx, \dzg@yt*\dzg@sy)
    node[above] {\dzg@ssQname};
  \ifdzg@ssticks
    \foreach \p [count=\k] in \dzg@ssP {%
      \ifx\dzg@ssPlabels\empty\def\dzg@t{$\p$}\else
        \foreach \l [count=\j] in \dzg@ssPlabels {\ifnum\j=\k\xdef\dzg@t{$\l$}\fi}\fi
      \node[font=\small, below] at (\p*\dzg@sx, \dzg@yb*\dzg@sy) {\dzg@t};}
    \foreach \q [count=\k] in \dzg@ssQ {%
      \ifx\dzg@ssQlabels\empty\def\dzg@t{$\q$}\else
        \foreach \l [count=\j] in \dzg@ssQlabels {\ifnum\j=\k\xdef\dzg@t{$\l$}\fi}\fi
      \node[font=\small, left] at (\dzg@xl*\dzg@sx, \q*\dzg@sy) {\dzg@t};}
  \fi}}}

\newcommand{\ssentry}[4]{\node (#1-#2-#3) at (#1-#2-#3) {#4};}
\newcommand{\ssdifferential}[5][]{%
  \begingroup
  \dzg@sslabeltrue
  \tikzset{#1}%
  \edef\dzg@c{\csname dzg@ssconv@#2\endcsname}\def\dzg@hom{hom}%
  \ifx\dzg@c\dzg@hom
    \pgfmathtruncatemacro\dzg@tp{#4-#3}\pgfmathtruncatemacro\dzg@tq{#5+#3-1}%
    \def\dzg@lab{$d^{#3}_{#4,#5}$}%
  \else
    \pgfmathtruncatemacro\dzg@tp{#4+#3}\pgfmathtruncatemacro\dzg@tq{#5-#3+1}%
    \def\dzg@lab{$d_{#3}^{#4,#5}$}%
  \fi
  \ifdzg@sslabel
    \draw[map arrow, shorten <=1pt, shorten >=1pt, #1] (#2-#4-#5) -- (#2-\dzg@tp-\dzg@tq)
      node[midway, fill=white, inner sep=1pt, font=\small] {\dzg@lab};
  \else
    \draw[map arrow, shorten <=1pt, shorten >=1pt, #1] (#2-#4-#5) -- (#2-\dzg@tp-\dzg@tq);
  \fi
  \endgroup}

% =============================================================================
% Hasse diagrams of serialized posets: constructors
% =============================================================================
% \posetfromjson[<keys>]{<file.json>}: the Hasse diagram of the poset in a
% networkx node-link JSON file, which Sage writes for any finite poset. The
% layout (lua/dzg-poset.lua) grades the poset, lets Graphviz dot order and
% space the elements of each grade, and places every element at the height
% of its grade. Each element is a node named by its id (so ids must be TikZ
% node names), an `element width` x `element height` box holding its label;
% \posetelements lists the ids, so a figure can place arbitrary content (a
% pic, a scope) in every box after the layout:
%   \posetfromjson[element width=2cm, element height=1.5cm]{m2.json}
%   \foreach \e in \posetelements \pic at (\e) {object=moduli/m2-graph-\e};
% Keys:
%   poset grading=<method>   used when the elements carry no `grade`:
%                            balanced (default), longest chain from bottom,
%                            longest chain to top, minimum total cover length
%                            (lua/dzg-poset.lua defines and compares them)
%   grade distance=<length>  the vertical distance per unit of grade
%   element width, element height   the box every element's content fits in
%   element sep=<length>     the least horizontal gap between two boxes
% Other keys apply to the scope holding the diagram.
% LuaLaTeX with --shell-escape only: the layout runs Graphviz dot.

\def\poset@grading{balanced}
\def\poset@distance{1cm}
\def\posetelementwidth{1cm}
\def\posetelementheight{0.6cm}
\def\poset@sep{3mm}
\tikzset{
  poset grading/.store in=\poset@grading,
  grade distance/.store in=\poset@distance,
  element width/.store in=\posetelementwidth,
  element height/.store in=\posetelementheight,
  element sep/.store in=\poset@sep,
  poset element/.style={shape=rectangle, inner sep=0pt,
    minimum width=\posetelementwidth, minimum height=\posetelementheight},
  poset cover/.style={draw, line width=0.6pt, -},
}
\ifdefined\directlua
  \directlua{dzgposet = require("dzg-poset")}
  \newcommand\posetfromjson[2][]{%
    \begingroup
    \tikzset{#1}%
    \pgfmathsetmacro\poset@w{(\posetelementwidth)/1cm}%
    \pgfmathsetmacro\poset@h{(\posetelementheight)/1cm}%
    \pgfmathsetmacro\poset@s{(\poset@sep)/1cm}%
    \pgfmathsetmacro\poset@d{(\poset@distance)/1cm}%
    \directlua{dzgposet.emit("\luaescapestring{#2}", "\luaescapestring{\unexpanded{#1}}",
      "\luaescapestring{\poset@grading}", \poset@w, \poset@h, \poset@s, \poset@d)}%
    \endgroup}
\else
  \newcommand\posetfromjson[2][]{\PackageError{dzg}{\string\posetfromjson\space
    needs LuaLaTeX}{The layout runs Graphviz dot from Lua; render the figure
    with lualatex --shell-escape.}}
\fi

